%%%%%%%%%%%%%%%%%%%%%%%%%%%%%%%%%%%%%%%%%%%%%%%%%%%%%%%%%%%%%%%%%%%%%%%%%%%
\chapter*{Glossary}
\addcontentsline{toc}{chapter}{Glossary}
%%%%%%%%%%%%%%%%%%%%%%%%%%%%%%%%%%%%%%%%%%%%%%%%%%%%%%%%%%%%%%%%%%%%%%%%%%%

Definitions for terms used throughout this book. Terms in \textbf{bold} within a definition
are themselves defined in this glossary.

\begin{description}

\item[\textbf{ANN (Approximate Nearest Neighbor)}]
A family of algorithms that find vectors close to a query vector without checking every vector
in the index. Trades a small amount of recall for large speedups over exact search. Common
algorithms: \textbf{HNSW}, \textbf{IVFFlat}.

\item[\textbf{Asymmetric Distance Computation (ADC)}]
A technique used with \textbf{Product Quantization} where the query vector is kept in full
precision but stored document vectors are quantized. The query is compared against lookup
tables of centroid distances.

\item[\textbf{Bi-encoder}]
A retrieval architecture where query and document are encoded independently by the same or
compatible models, producing separate vectors. Fast because document vectors can be precomputed.
Contrast with \textbf{cross-encoder}.

\item[\textbf{BM25 (Best Match 25)}]
A full-text ranking algorithm based on term frequency, inverse document frequency, and
document length normalization. The standard algorithm behind Elasticsearch, Lucene, Solr, and
Tantivy. Excels at exact keyword matching and rare term retrieval.

\item[\textbf{BPE (Byte Pair Encoding)}]
A \textbf{tokenization} algorithm that iteratively merges the most frequent adjacent byte
pairs into a single token. Used by GPT-2, GPT-3, GPT-4, LLaMA, Qwen, and most modern LLMs.

\item[\textbf{Chunking}]
Splitting a document into smaller pieces before embedding and indexing. Necessary because
embedding models have token limits and because smaller chunks improve retrieval precision.

\item[\textbf{ColBERT}]
A retrieval model that keeps one vector per token rather than collapsing to a single vector.
Similarity is computed via \textbf{MaxSim}: each query token finds its best-matching document
token. Higher recall on complex queries; larger index.

\item[\textbf{Contrastive learning}]
A training objective where the model learns to bring similar examples close together and push
dissimilar examples apart in vector space. Uses (anchor, positive, negative) triplets.

\item[\textbf{Cosine similarity}]
A distance metric measuring the angle between two vectors, regardless of magnitude. Values
range from $-1$ (opposite) to $1$ (same direction). The standard metric for most embedding
models. Equivalent to \textbf{dot product} on unit-normalized vectors.

\item[\textbf{Cross-encoder}]
A model that reads query and document together as a single input, producing a single relevance
score. More accurate than \textbf{bi-encoders} because it reasons over query-document
interaction. Slower because it cannot precompute. Used as a \textbf{reranker}.

\item[\textbf{Dense retrieval}]
Retrieval using a single continuous vector per text. Contrast with \textbf{sparse retrieval}
and \textbf{ColBERT}.

\item[\textbf{Embedding}]
A fixed-size vector of floating-point numbers representing the semantic content of a text.
Produced by an embedding model. Similar texts produce similar (nearby) vectors.

\item[\textbf{GGUF}]
A binary file format for quantized model weights, used by llama.cpp and llama-server. Enables
running large models on consumer hardware by reducing precision (e.g., Q4\_K\_M, Q8\_0).

\item[\textbf{Hard negative}]
A training example that appears relevant to the query but is actually not the correct match.
Hard negatives force the embedding model to learn fine-grained distinctions.

\item[\textbf{HNSW (Hierarchical Navigable Small World)}]
The dominant \textbf{ANN} algorithm. Builds a multi-layer graph of vectors; upper layers are
sparse for long-range navigation, lower layers are dense for local search. Key parameters:
M, ef\_construction, ef.

\item[\textbf{HyDE (Hypothetical Document Embedding)}]
A query transformation technique: ask an LLM to write a hypothetical document that would
answer the question, then embed that document as the query vector. Bridges the vocabulary gap
between short queries and long documents.

\item[\textbf{Hybrid search}]
Combining multiple retrieval modalities (typically \textbf{BM25} and \textbf{dense retrieval})
and fusing the results with \textbf{RRF} or other merging strategies. Consistently outperforms
either modality alone.

\item[\textbf{L2 distance (Euclidean distance)}]
Straight-line distance between two vectors: $\sqrt{\sum_i (a_i - b_i)^2}$. Smaller = more
similar. Equivalent to cosine distance on unit-normalized vectors.

\item[\textbf{Late chunking}]
Embedding the full document first to preserve context, then extracting chunk-level vectors
from the token-level outputs. Avoids boundary artifacts. Requires token-level model outputs.

\item[\textbf{MemoryCard}]
A typed fact record in memvid's memory system. Types: Fact, Preference, Event, Relationship.
Built on top of \textbf{SPO} triplets.

\item[\textbf{MTEB (Massive Text Embedding Benchmark)}]
The primary benchmark suite for evaluating embedding models across multiple task types
(retrieval, classification, clustering, reranking, STS). High MTEB scores may reflect
benchmark proximity rather than out-of-distribution generalization.

\item[\textbf{nDCG (Normalized Discounted Cumulative Gain)}]
The primary metric in most retrieval benchmarks. Measures ranking quality: highly relevant
documents appearing earlier receive higher scores. Values range from 0 to 1; 1 = perfect
ranking. \texttt{nDCG@10} = nDCG considering only the top 10 results.

\item[\textbf{Pooling}]
The step that collapses per-token vectors from a transformer encoder into a single vector.
Common strategies: mean pooling, CLS pooling, last-token pooling. The pooling strategy is
part of the model contract and must match how the model was trained.

\item[\textbf{PQ (Product Quantization)}]
A vector compression technique that splits each vector into M subspaces and quantizes each
subspace to one of k centroids. Reduces storage by 4--64$\times$ at the cost of some recall.

\item[\textbf{Query instruction}]
A task-specific prefix added to the query before embedding. Improves retrieval quality by
1--5\% on models trained with instructions (Qwen3, nomic-embed-text).

\item[\textbf{RAG (Retrieval-Augmented Generation)}]
A pattern combining retrieval (finding relevant documents) with generation (producing an
answer). The LLM generates answers conditioned on retrieved context rather than relying
solely on parametric memory.

\item[\textbf{Recall@k}]
The fraction of all relevant documents that appear in the top-$k$ retrieved results.
Contrast with \textbf{Precision@k}.

\item[\textbf{Reranker}]
A \textbf{cross-encoder} model used as a second pass over candidates returned by first-stage
retrieval. Produces more accurate relevance scores. Applied to a small candidate pool
(20--100 items).

\item[\textbf{RRF (Reciprocal Rank Fusion)}]
A method for merging multiple ranked result lists by position rather than score. Each result
scores $\sum 1 / (k + \text{rank})$. Requires no score normalization.

\item[\textbf{RTEB (Retrieval Text Embedding Benchmark)}]
A benchmark designed to test embedding models on retrieval-specific tasks with held-out
evaluation datasets not seen during training. Intended to detect benchmark overfitting.

\item[\textbf{SPLADE}]
A sparse retrieval model that produces learned sparse vectors (one weight per vocabulary
term) rather than dense vectors. Complementary to dense retrieval in hybrid search.

\item[\textbf{SPO (Subject-Predicate-Object)}]
A structured representation of a fact as a triple. Used to build entity graphs in agent
memory systems.

\item[\textbf{Tokenization}]
The process of splitting text into tokens---the basic units a model processes. Common
algorithms: \textbf{BPE}, WordPiece, SentencePiece.

\item[\textbf{Transformer}]
The neural network architecture underlying all modern embedding and generation models.
Key components: multi-head self-attention, feed-forward layers, positional encoding.

\end{description}
